
\documentclass[aspectratio=169,10pt]{beamer}

% ==========================================
% PAQUETES
% ==========================================
\usepackage[utf8]{inputenc}
\usepackage[T1]{fontenc}
\usepackage[spanish]{babel}
\usepackage{lmodern}
\usepackage{amsmath,amssymb}
\usepackage{tikz}
\usepackage{xcolor}
\usepackage{array}

% TikZ antes de begin document
\usetikzlibrary{
 arrows.meta,
 calc,
 positioning,
 patterns
}

% ==========================================
% COLORES Y DISEÑO
% ==========================================
\definecolor{azul}{RGB}{0,100,185}
\definecolor{azuloscuro}{RGB}{15,60,115}
\definecolor{verde}{RGB}{0,145,95}
\definecolor{morado}{RGB}{115,55,175}
\definecolor{naranja}{RGB}{240,105,15}
\definecolor{celeste}{RGB}{230,243,255}

\usetheme{Madrid}

\setbeamercolor{structure}{fg=azuloscuro}
\setbeamercolor{frametitle}{
 bg=azuloscuro,
 fg=white
}
\setbeamercolor{block title}{
 bg=azul,
 fg=white
}
\setbeamercolor{block body}{
 bg=celeste,
 fg=black
}
\setbeamercolor{block title example}{
 bg=verde,
 fg=white
}
\setbeamercolor{block body example}{
 bg=green!5,
 fg=black
}

\setbeamertemplate{navigation symbols}{}
\setbeamertemplate{footline}[frame number]

% ==========================================
% COMANDOS
% ==========================================
\newcommand{\resultado}[1]{
\[
\boxed{\color{morado}#1}
\]
}

\newcommand{\datos}[1]{
\textbf{Datos:}\quad #1
\par\medskip
}

\newcommand{\formula}[1]{
\textbf{Fórmula:}
\[
#1
\]
}

\newcommand{\sustitucion}[1]{
\textbf{Sustitución:}
\[
#1
\]
}

% Gráfica de posición-tiempo reutilizable
% Argumento: lista de coordenadas
\newcommand{\grafpos}[1]{
\begin{tikzpicture}[
 x=0.11cm,
 y=0.075cm,
 >=Stealth
]
\draw[->,thick]
(0,0)--(57,0)
node[right] {$t$ (s)};

\draw[->,thick]
(0,0)--(0,46)
node[above] {$d$ (m)};

\foreach \x in {10,20,30,40,50}{
 \draw (\x,0)--(\x,-1);
 \node[below,font=\scriptsize]
 at (\x,-1) {\x};
}

\foreach \y in {10,20,30,40}{
 \draw (0,\y)--(-1,\y);
 \node[left,font=\scriptsize]
 at (-1,\y) {\y};
}

\draw[
 very thick,
 azul,
 mark=*
] plot coordinates {#1};

\end{tikzpicture}
}

% ==========================================
% PORTADA
% ==========================================
\title{Gráficas del movimiento rectilíneo}

\subtitle{
Distancia, desplazamiento, rapidez,
velocidad, pendiente y áreas
}

\author{Clases de Matemática Paraíso}

\date{Física -- Educación Diversificada}

\begin{document}

% ==========================================
% DIAPOSITIVA 1
% ==========================================
\begin{frame}
\titlepage
\end{frame}

% ==========================================
% DIAPOSITIVA 2
% ==========================================
\begin{frame}{Objetivos de aprendizaje}

Al finalizar esta unidad, el estudiante podrá:

\begin{enumerate}
\item Interpretar gráficas de movimiento.

\item Calcular rapidez y velocidad
mediante la pendiente.

\item Determinar distancia recorrida
y desplazamiento.

\item Identificar intervalos de reposo.

\item Calcular distancias mediante
áreas bajo gráficas rapidez--tiempo.

\item Interpretar gráficas
velocidad--tiempo.
\end{enumerate}

\end{frame}

% ==========================================
% DIAPOSITIVA 3
% ==========================================
\section{Conceptos fundamentales}

\begin{frame}{Distancia y desplazamiento}

\begin{block}{Distancia recorrida}
Es la longitud total del camino recorrido.

Es una magnitud escalar y nunca es negativa.
\end{block}

\begin{block}{Desplazamiento}
Es el cambio de posición de un cuerpo.

\[
\Delta d=d_f-d_i
\]

Puede ser positivo, negativo o cero.
\end{block}

\begin{exampleblock}{Ejemplo}
Un estudiante camina $30\,\mathrm{m}$
hacia el este y regresa $10\,\mathrm{m}$.

\[
\text{Distancia}=30+10=40\,\mathrm{m}
\]

\[
\Delta d=30-10=20\,\mathrm{m}
\]
\end{exampleblock}

\end{frame}

% ==========================================
% DIAPOSITIVA 4
% ==========================================
\begin{frame}{Rapidez y velocidad}

\begin{block}{Rapidez media}
Relaciona la distancia total
con el tiempo transcurrido.

\[
v=\frac{\text{distancia total}}{t}
\]
\end{block}

\begin{block}{Velocidad media}
Relaciona el desplazamiento
con el tiempo transcurrido.

\[
v_m=\frac{\Delta d}{\Delta t}
\]
\end{block}

\begin{itemize}
\item Rapidez: no indica dirección.
\item Velocidad: incluye sentido del movimiento.
\end{itemize}

\end{frame}

% ==========================================
% DIAPOSITIVA 5
% ==========================================
\section{Gráficas posición-tiempo}

\begin{frame}{Gráficas posición--tiempo}

Una gráfica posición--tiempo representa
cómo cambia la posición de un cuerpo.

\begin{center}
\grafpos{
(0,0)
(10,20)
(30,30)
(40,30)
(50,10)
}
\end{center}

\begin{block}{Interpretación}
\begin{itemize}
\item Recta ascendente: movimiento positivo.
\item Recta horizontal: reposo.
\item Recta descendente: movimiento negativo.
\end{itemize}
\end{block}

\small
Aquí $d$ representa la posición
respecto al origen, no la distancia
total acumulada.

\end{frame}

% ==========================================
% DIAPOSITIVA 6
% ==========================================
\begin{frame}{La pendiente de una gráfica}

\begin{block}{Concepto fundamental}

La pendiente de una gráfica
posición--tiempo representa la velocidad.

\[
\boxed{
v=\frac{d_f-d_i}{t_f-t_i}
}
\]

\end{block}

\begin{center}
\begin{tikzpicture}[
 x=0.22cm,
 y=0.10cm,
 >=Stealth
]

\draw[->,thick]
(0,0)--(25,0)
node[right] {$t$ (s)};

\draw[->,thick]
(0,0)--(0,35)
node[above] {$d$ (m)};

\draw[very thick,azul]
(0,0)--(20,30);

\draw[dashed,gray]
(20,0)--(20,30)--(0,30);

\node[below] at (20,0) {$20$};
\node[left] at (0,30) {$30$};

\end{tikzpicture}
\end{center}

\[
v=\frac{30-0}{20-0}
=1,5\,\mathrm{m/s}
\]

\end{frame}

% ==========================================
% DIAPOSITIVA 7
% ==========================================
\begin{frame}{Ejemplo resuelto: posición--tiempo}

\begin{columns}[T]

\column{0.53\textwidth}

\begin{center}
\grafpos{
(0,0)
(10,20)
(30,30)
(40,30)
(50,10)
}
\end{center}

\column{0.47\textwidth}

\begin{block}{Determine}
\begin{enumerate}
\item Velocidad en cada tramo.
\item Distancia total.
\item Desplazamiento.
\item Intervalo de reposo.
\end{enumerate}
\end{block}

\end{columns}

\end{frame}

% ==========================================
% DIAPOSITIVA 8
% ==========================================
\begin{frame}{Solución: velocidades por intervalo}

\small

\textbf{Tramo I: de 0 a 10 s}

\[
v_1=\frac{20-0}{10-0}
=2\,\mathrm{m/s}
\]

\textbf{Tramo II: de 10 a 30 s}

\[
v_2=\frac{30-20}{30-10}
=0,5\,\mathrm{m/s}
\]

\textbf{Tramo III: de 30 a 40 s}

\[
v_3=\frac{30-30}{40-30}
=0\,\mathrm{m/s}
\]

\textbf{Tramo IV: de 40 a 50 s}

\[
v_4=\frac{10-30}{50-40}
=-2\,\mathrm{m/s}
\]

\end{frame}

% ==========================================
% DIAPOSITIVA 9
% ==========================================
\begin{frame}{Solución: distancia y desplazamiento}

\textbf{Distancia total}

\[
D=|20-0|+|30-20|
+|30-30|+|10-30|
\]

\[
D=20+10+0+20
\]

\resultado{D=50\,\mathrm{m}}

\textbf{Desplazamiento}

\[
\Delta d=d_f-d_i
\]

\[
\Delta d=10-0
\]

\resultado{\Delta d=10\,\mathrm{m}}

\begin{exampleblock}{Reposo}
El móvil permanece en reposo
entre $30\,\mathrm{s}$ y $40\,\mathrm{s}$.
\end{exampleblock}

\end{frame}

% ==========================================
% DIAPOSITIVA 10
% ==========================================
\begin{frame}{Práctica: posición--tiempo}

\begin{center}
\grafpos{
(0,10)
(10,30)
(20,30)
(30,0)
(40,20)
}
\end{center}

Determine:

\begin{enumerate}
\item Velocidad de cada intervalo.
\item Distancia total recorrida.
\item Desplazamiento.
\item Intervalo de reposo.
\item Tramo de mayor rapidez.
\end{enumerate}

\end{frame}

% ==========================================
% DIAPOSITIVA 11
% ==========================================
\section{Gráficas rapidez-tiempo}

\begin{frame}{Gráficas rapidez--tiempo}

\begin{block}{Concepto fundamental}

En una gráfica rapidez--tiempo,
el área bajo la curva representa
la distancia recorrida.

\[
\boxed{D=\text{Área bajo la curva}}
\]

\end{block}

\begin{center}
\begin{tikzpicture}[
 x=0.8cm,
 y=0.35cm,
 >=Stealth
]

\draw[->,thick]
(0,0)--(9,0)
node[right] {$t$ (s)};

\draw[->,thick]
(0,0)--(0,12)
node[above] {$v$ (m/s)};

\fill[blue!15]
(0,0)--(2,6)--(4,6)
--(6,10)--(8,0)--cycle;

\draw[very thick,azul]
(0,0)--(2,6)--(4,6)
--(6,10)--(8,0);

\foreach \x in {2,4,6,8}{
\node[below] at (\x,0) {\x};
}

\foreach \y in {2,4,6,8,10}{
\node[left] at (0,\y) {\y};
}

\end{tikzpicture}
\end{center}

\end{frame}

% ==========================================
% DIAPOSITIVA 12
% ==========================================
\begin{frame}{Figuras geométricas y áreas}

Para calcular la distancia,
utilizamos las fórmulas de áreas.

\begin{block}{Rectángulo}
\[
A=b\cdot h
\]
\end{block}

\begin{block}{Triángulo}
\[
A=\frac{b\cdot h}{2}
\]
\end{block}

\begin{block}{Trapecio}
\[
A=\frac{(B+b)h}{2}
\]
\end{block}

En una gráfica rapidez--tiempo:

\[
\mathrm{s}\cdot\frac{\mathrm{m}}{\mathrm{s}}
=\mathrm{m}
\]

\end{frame}

% ==========================================
% DIAPOSITIVA 13
% ==========================================
\begin{frame}{Ejemplo: distancia mediante áreas}

\begin{center}
\begin{tikzpicture}[
 x=0.85cm,
 y=0.38cm,
 >=Stealth
]

% Áreas
\fill[blue!17]
(0,0)--(2,6)--(2,0)--cycle;

\fill[green!18]
(2,0) rectangle (4,6);

\fill[orange!17]
(4,0)--(4,6)--(6,10)--(6,0)--cycle;

\fill[purple!15]
(6,0)--(6,10)--(8,0)--cycle;

% Ejes
\draw[->,thick]
(0,0)--(9,0)
node[right] {$t$ (s)};

\draw[->,thick]
(0,0)--(0,12)
node[above] {$v$ (m/s)};

% Curva
\draw[very thick,azul]
(0,0)--(2,6)--(4,6)
--(6,10)--(8,0);

% Divisiones
\foreach \x in {2,4,6}{
\draw[dashed,gray]
(\x,0)--(\x,{ifthenelse(\x==6,10,6)});
}

% Etiquetas
\node at (1,1.5) {I};
\node at (3,3) {II};
\node at (5,3) {III};
\node at (6.7,3) {IV};

\foreach \x in {2,4,6,8}{
\node[below] at (\x,0) {\x};
}

\node[left] at (0,6) {6};
\node[left] at (0,10) {10};

\end{tikzpicture}
\end{center}

Determine la distancia recorrida
en cada intervalo y la distancia total.

\end{frame}

% ==========================================
% DIAPOSITIVA 14
% ==========================================
\begin{frame}{Solución: áreas I y II}

\textbf{Intervalo I}

Triángulo:

\[
A_1=\frac{b\cdot h}{2}
\]

\[
A_1=\frac{(2)(6)}{2}
\]

\resultado{A_1=6\,\mathrm{m}}

\textbf{Intervalo II}

Rectángulo:

\[
A_2=b\cdot h
\]

\[
A_2=(2)(6)
\]

\resultado{A_2=12\,\mathrm{m}}

\end{frame}

% ==========================================
% DIAPOSITIVA 15
% ==========================================
\begin{frame}{Solución: áreas III y IV}

\textbf{Intervalo III}

Trapecio:

\[
A_3=\frac{(B+b)h}{2}
\]

\[
A_3=\frac{(10+6)(2)}{2}
\]

\resultado{A_3=16\,\mathrm{m}}

\textbf{Intervalo IV}

Triángulo:

\[
A_4=\frac{(2)(10)}{2}
\]

\resultado{A_4=10\,\mathrm{m}}

\end{frame}

% ==========================================
% DIAPOSITIVA 16
% ==========================================
\begin{frame}{Distancia total recorrida}

La distancia total corresponde
a la suma de todas las áreas.

\[
D=A_1+A_2+A_3+A_4
\]

Sustituimos:

\[
D=6+12+16+10
\]

\resultado{D=44\,\mathrm{m}}

\begin{exampleblock}{Interpretación}

Durante los ocho segundos,
el móvil recorrió una distancia
total de $44\,\mathrm{m}$.

\end{exampleblock}

\end{frame}

% ==========================================
% DIAPOSITIVA 17
% ==========================================
\section{Gráficas velocidad-tiempo}

\begin{frame}{Gráficas velocidad--tiempo}

\begin{block}{Conceptos fundamentales}

La pendiente representa la aceleración.

\[
\boxed{
a=\frac{v_f-v_i}{t_f-t_i}
}
\]

El área con signo representa
el desplazamiento.

\[
\boxed{
\Delta d=\text{Área con signo}
}
\]

\end{block}

\begin{itemize}
\item Área sobre el eje: desplazamiento positivo.
\item Área bajo el eje: desplazamiento negativo.
\item La suma de áreas absolutas da
la distancia total recorrida.
\end{itemize}

\end{frame}


% ==========================================
% DIAPOSITIVA 18
% ==========================================
\begin{frame}{Ejemplo: velocidad negativa}

\begin{center}
\begin{tikzpicture}[
 x=0.65cm,
 y=0.25cm,
 >=Stealth,
 font=\footnotesize
]

% Ejes
\draw[->,thick]
(0,0)--(9,0)
node[right] {$t$ (s)};

\draw[->,thick]
(0,-4.5)--(0,6)
node[above] {$v$ (m/s)};

% Áreas
\fill[blue!20]
(0,0) rectangle (4,4);

\fill[orange!20]
(4,-3) rectangle (8,0);

% Velocidades
\draw[very thick,azul]
(0,4)--(4,4);

\draw[very thick,naranja]
(4,-3)--(8,-3);

% Cambio instantáneo
\draw[dashed,gray]
(4,4)--(4,-3);

% Valores
\node[left] at (0,4) {$4$};
\node[left] at (0,-3) {$-3$};

\node[below right] at (4,0) {$4$};
\node[below] at (8,0) {$8$};

\end{tikzpicture}
\end{center}

\vspace{-0.25cm}

{\footnotesize
En este ejemplo, el cambio de velocidad en
$t=4\,\mathrm{s}$ es idealizado e instantáneo.
}

\vspace{-0.15cm}

\begingroup
\setlength{\abovedisplayskip}{3pt}
\setlength{\belowdisplayskip}{3pt}

\small

\[
A_1=(4)(4)=16\,\mathrm{m}
\]

\[
A_2=(4)(-3)=-12\,\mathrm{m}
\]

\begin{center}
\resultado{\Delta d=16-12=4\,\mathrm{m}}
\end{center}

\[
D=16+12=28\,\mathrm{m}
\]

\endgroup

\end{frame}

% ==========================================
% DIAPOSITIVA 19
% ==========================================
\begin{frame}{Resumen de la unidad}

\begin{center}
\renewcommand{\arraystretch}{1.6}

\begin{tabular}{|p{3.3cm}|p{6.3cm}|}
\hline
\textbf{Gráfica}
&
\textbf{Interpretación}\\
\hline

Posición--tiempo
&
Pendiente: velocidad\\
\hline

Distancia acumulada--tiempo
&
Pendiente: rapidez\\
\hline

Rapidez--tiempo
&
Área: distancia recorrida\\
\hline

Velocidad--tiempo
&
Pendiente: aceleración\\
\hline

Velocidad--tiempo
&
Área con signo: desplazamiento\\
\hline

\end{tabular}
\end{center}

\vspace{0.4cm}

\begin{exampleblock}{Recuerde}

Antes de resolver un ejercicio,
identifique las magnitudes de los ejes,
sus unidades y los intervalos.

\end{exampleblock}

\end{frame}

\end{document}
